% Einheit 2 – Wachstumsfaktor und Wachstumstabelle
\setcounter{aufgabe}{20}
\hypertarget{e2}{}\einheitenkopf*{Einheit 2 von 5 · Wachstumsfaktor und Wachstumstabelle}
\zweigzeile{Hier lernst du, aus dem Prozentsatz den Wachstumsfaktor zu bilden und eine Wachstumstabelle mit ihm fortzuschreiben \,·\, neu in diesem Jahr \,·\, P10 oft \,·\, baut auf: Einheit 1, Prozentwert und Prozentsatz, Zinseszins, Quotient zweier Werte}

\begin{aufgabe}{Ich erkenne am Faktor, ob ein Wert zunimmt oder abnimmt. Ein Wert wird jedes Jahr mit der angegebenen Zahl malgenommen. Kreuze an. Du musst nichts rechnen.}
\begin{teilezwei}
\tz mal $1{,}15$: \ \kreuz{nimmt zu}\kreuz{nimmt ab} & \tz mal $0{,}7$: \ \kreuz{nimmt zu}\kreuz{nimmt ab} \\
\tz mal $2$: \ \kreuz{nimmt zu}\kreuz{nimmt ab} & \tz mal $0{,}98$: \ \kreuz{nimmt zu}\kreuz{nimmt ab} \\
\tz mal $1{,}005$: \ \kreuz{nimmt zu}\kreuz{nimmt ab} & \\
\end{teilezwei}
\end{aufgabe}

\begin{aufgabe}{Ich kann den Wachstumsfaktor bestimmen. Ein Wert wächst jedes Jahr um den angegebenen Prozentsatz. Gib den Wachstumsfaktor $q$ an.}
\beispiel{\text{Zunahme um } 11\,\% &\;\rightarrow\; q = 1 + \tfrac{11}{100} = 1{,}11}
\begin{teilezwei}
\tz Zunahme um $10\,\%$: \feld{q} & \tz Zunahme um $15\,\%$: \feld{q} \\
\tz Zunahme um $2\,\%$: \feld{q} & \tz Zunahme um $35\,\%$: \feld{q} \\
\end{teilezwei}
Bei einer Abnahme ist $q = 1 - \tfrac{p}{100}$; $p$ ist der Prozentsatz.
\begin{teile}
\teil Abnahme um $40\,\%$: \feld{q}
\end{teile}
Gib an, um wie viel Prozent der Wert jedes Jahr zu- oder abnimmt.
\begin{teile}
\teil $q = 0{,}93$: \quad \kreuz{Zunahme}\kreuz{Abnahme} um \leerfeld[\%]
\end{teile}
Berechne den Wachstumsfaktor als Quotient: späterer Wert geteilt durch früheren Wert.
\begin{teile}
\teil Ein Wert wächst in einem Jahr von $250$ auf $290$. \quad \feld{q}
\teil Werte nach 0, 1, 2, 3 Jahren: $400;\ 460;\ 529;\ 608{,}35$. Bilde alle Quotienten und schreibe das Ergebnis in einem Satz.
\teil Ein Wert wächst jedes Jahr um denselben Prozentsatz. Am Anfang sind es $1\,600$, nach zwei Jahren $1\,936$. Bestimme den Prozentsatz. \quad \leerfeld[\%]
\teil Die Tabelle zeigt die Zahl der Hefezellen in einer Probe. Weise nach, dass sich die Zahl jede Stunde mit demselben Faktor vervielfacht, und gib den Faktor an. (P10 2025)
\end{teile}
\sachtabelle{lcccc}{Zeit in Stunden & 0 & 1 & 2 & 3}{Hefezellen & 256 & 320 & 400 & 500}
\end{aufgabe}

\begin{aufgabe}{Ich kann eine Wachstumstabelle mit dem Wachstumsfaktor fortschreiben. Berechne den nächsten Wert.}
\beispiel{q = 1{,}1:\quad 300 &\;\rightarrow\; 300 \cdot 1{,}1 = 330}
\begin{teilezwei}
\tz $70$, \ $q = 2$: \ \leerfeld & \tz $500$, \ $q = 1{,}1$: \ \leerfeld \\
\tz $800$, \ $q = 1{,}05$: \ \leerfeld & \tz $250$, \ $q = 1{,}02$: \ \leerfeld \\
\end{teilezwei}
Ergänze die Tabelle.
\begin{teile}
\teil Ein Bestand beginnt mit $400$ Tieren und wächst jedes Jahr um $10\,\%$.
\end{teile}
\sachtabelle{lccc}{Jahr & 0 & 1 & 2}{Tiere & \leerzelle & \leerzelle & \leerzelle}
\begin{teile}
\teil Der Faktor ist $q = 1{,}1$. Fülle die Lücke.
\end{teile}
\sachtabelle{lccc}{Jahr & 0 & 1 & 2}{Wert & 1\,500 & \leerzelle & 1\,815}
\begin{teile}
\teil $3\,000\,€$ wachsen jedes Jahr um $5\,\%$. Berechne den Wert nach $3$ Jahren in einem Schritt mit der Potenz des Faktors. \quad \leerfeld[€]
\end{teile}
\end{aufgabe}

\begin{aufgabe}{Ich kann eine Wachstumstabelle mit dem Wachstumsfaktor fortschreiben – weiter.}
\begin{teile}
\teil Im Jahr 3 leben in einem Park $1\,331$ Kaninchen, der Faktor ist $q = 1{,}1$. Wie viele waren es im Jahr 2? \quad \leerfeld
\teil Der Faktor ist $q = 1{,}1$. Bestimme die fehlende Jahreszahl und prüfe mit dem Faktor.
\end{teile}
\sachtabelle{lccc}{Jahr & 0 & 1 & \leerzelle}{Wert & 500 & 550 & 732{,}05}
\begin{teile}
\teil Der Faktor ist $q = 0{,}8$. Ergänze die Tabelle.
\end{teile}
\sachtabelle{lcccc}{Jahr & 0 & 1 & 2 & 3}{Wert & 750 & 600 & \leerzelle & \leerzelle}
\begin{teile}
\teil Ein Streaming-Abo kostet im Jahr 2025 insgesamt $144\,€$. Der Preis steigt jedes Jahr um $3{,}5\,\%$. Ergänze die fehlenden Werte der Tabelle. (P10 2026)
\end{teile}
\sachtabelle{lcccc}{Jahr & 2025 & 2026 & 2027 & 2028}{Kosten in € & \leerzelle & 149{,}04 & \leerzelle & 159{,}66}
\end{aufgabe}

\begin{aufgabe}{Ich finde den Fehler beim Fortschreiben einer Wachstumstabelle. Ein Motorrad kostet neu $7\,200\,€$ und verliert jedes Jahr $15\,\%$ an Wert. Paul schreibt die Tabelle so fort:}
\sachtabelle{lcccc}{Jahr & 0 & 1 & 2 & 3}{Wert in € & 7\,200 & 6\,120 & 5\,040 & 3\,960}
\begin{teile}
\teil Finde den Fehler, benenne ihn und korrigiere die Tabelle.
\teil Ein Roller kostet neu $3\,600\,€$ und verliert jedes Jahr $10\,\%$ an Wert. Schreibe die Tabelle bis zum Jahr 3 fort.
\end{teile}
\sachtabelle{lcccc}{Jahr & 0 & 1 & 2 & 3}{Wert in € & 3\,600 & \leerzelle & \leerzelle & \leerzelle}
Lena soll $2\,000\,€$ um $9\,\%$ erhöhen. Sie rechnet: $2\,000 \cdot 0{,}09 = 180$, also $180\,€$.
\begin{teile}
\teil Finde den Fehler, benenne ihn und korrigiere die Rechnung.
\teil Erhöhe selbst $2\,500\,€$ um $7\,\%$ mit dem Wachstumsfaktor. \quad \leerfeld[€]
\end{teile}
\end{aufgabe}

\begin{aufgabe}{Ich kann begründen, warum der Quotient den Wachstumsfaktor zeigt. Entscheide ohne Rechnung und begründe.}
\begin{teile}
\teil Warum bedeutet der Wachstumsfaktor $0{,}9$ eine Abnahme um $10\,\%$?
\teil Warum zeigt der Quotient zweier Nachbarwerte einer Tabelle den Wachstumsfaktor, die Differenz aber nicht?
\teil Tom sagt: „Wenn ein Preis erst um $50\,\%$ steigt und dann um $50\,\%$ sinkt, ist er wieder so hoch wie am Anfang.“ Hat Tom recht?
\end{teile}
\end{aufgabe}

\begin{aufgabe}{Ich kann mit einer Wachstumstabelle eine Entscheidung treffen. Eine Gemeinde hat $12\,000$ Einwohner. Die Zahl wächst jedes Jahr um $2\,\%$. Wohnen dort mehr als $12\,700$ Menschen, braucht die Grundschule einen Anbau.}
\begin{teile}
\teil Lege eine Tabelle für die nächsten drei Jahre an. Runde auf ganze Einwohner.
\teil Braucht die Grundschule in den nächsten drei Jahren einen Anbau? Wenn ja: ab welchem Jahr?
\end{teile}
\end{aufgabe}
